\documentclass[aspectratio=169,10pt]{beamer}

%% Use the official UTB theme (Trường Đại học Tây Bắc)
%% Options:
%%   color       = green | navy | blue
%%   font        = rokkit | lato | arial | times | cm
%%   titlefont   = same | avantgarde | barberfill
%%   footerlogo  = mark | full | wordmark | utb | none
%%   bar         = thin | thick
%%   titlepage   = cover | solid | light | dark | welcome
%%   coverbg     = solid | light
%%   sectionpage = none | panel
%%   numbering   = none | plain | fraction
%%   progressbar = on | off
%%   footer      = <custom footer text>
%%   footerleft  = <custom left text>
%%   logolayout  = 1 | 2 | 3 | 4a | 4b
\usetheme[
  color=green,
  font=lato,
  titlefont=avantgarde,
  footerlogo=mark,
  bar=thin,
  titlepage=cover,
  sectionpage=panel,
  numbering=fraction,
  progressbar=on,
  dateaddress=on,
  address={Sơn La},
  logolayout=2,
  logo2image=nst-shield-utb.png,
  logo2vn={KHOA KHOA HỌC TỰ NHIÊN - CÔNG NGHỆ},
  logo2en={Faculty of Natural Sciences and Technology},
  logo2color=utbnavy,
  logo2layout=1,
  headerlogo=utb-shield.png,
  framebg=on,
  footermeta=on
]{UTB}

%% ------------------------------------------------------------- Modular Academic & UX Setup
%% ==============================================================================
%% UTB Beamer Theme - Setup & Academic Extension Package (utb-setup.tex)
%% Trường Đại học Tây Bắc (UTB / Tay Bac University)
%% Modular setup file containing academic packages, UX/UI extensions,
%% math operators, list customizations, and cross-references.
%% ==============================================================================

%% ------------------------------------------------------------- 1. Layout & Tables
\RequirePackage{wrapfig}
\RequirePackage{booktabs}
\RequirePackage{tabularx}
\RequirePackage{multirow, multicol}
\RequirePackage{array}
\renewcommand{\tabularxcolumn}[1]{m{#1}}
\RequirePackage{caption}
\captionsetup{font=scriptsize, labelfont={bf, color=utbmain}, labelsep=colon}

%% TikZ libraries for graphics & markers:
\usetikzlibrary{calc, decorations.markings, arrows.meta}

%% ------------------------------------------------------------- 2. List Customization (enumitem & tasks)
\RequirePackage{enumitem}
\setlist[enumerate,1]{label=\arabic*.}
\setlist[enumerate,2]{label=\alph*.}
\setlist[enumerate,3]{label=\roman*)}
\setlist[itemize,1]{label=\usebeamerfont*{itemize item}\usebeamercolor[fg]{itemize item}\raise0.2ex\hbox{\vrule height 0.8ex width 0.8ex}}
\setlist[itemize,2]{label=\usebeamerfont*{itemize subitem}\usebeamercolor[fg]{itemize subitem}\raise0.2ex\hbox{\vrule height 0.6ex width 0.6ex}}
\setlist[itemize,3]{label=\raisebox{0.2ex}{\tiny\textcolor{utbaccent}{$\bullet$}}}
\setlist[itemize,4]{label=\raisebox{0.2ex}{\tiny\textcolor{utbmain}{$-$}}}

%% Square number box badge matching UTB Brand identity:
\DeclareRobustCommand{\utbitemsq}[1]{%
  \raisebox{0.05ex}{\colorbox{utbmain}{\makebox[1.05em]{\textcolor{white}{\bfseries\scriptsize#1}}}}%
}
\let\hustitemsq\utbitemsq

%% Dedicated square-numbered list environment:
\newlist{utbsqlist}{enumerate}{2}
\setlist[utbsqlist,1]{label=\utbitemsq{\arabic*}, ref=\arabic*, leftmargin=2.2em, labelsep=0.6em, itemsep=4pt, topsep=2pt}
\setlist[utbsqlist,2]{label=\utbitemsq{\alph*}, ref=\alph*, leftmargin=2.2em, labelsep=0.6em, itemsep=3pt, topsep=2pt}

\newenvironment{hustsqlist}{\begin{utbsqlist}}{\end{utbsqlist}}

%% Tasks package for multi-column horizontal listings:
\RequirePackage{tasks}
\settasks{
  before-skip = 2pt,
  after-skip = 2pt,
  after-item-skip = 3pt,
  label-width = 5pt,
  item-indent = 9pt,
  label-offset = 2pt,
  column-sep = 1em,
  ref = \arabic*,
  label = \raisebox{0.15ex}{\textcolor{utbmain}{\rule{0.75ex}{0.75ex}}}
}

%% ------------------------------------------------------------- 3. Mathematics & Operators
\RequirePackage{amsmath, mathtools}
\RequirePackage{derivative}
\RequirePackage{nicematrix, nicefrac}
\RequirePackage{siunitx}

\ifdefined\symbfit
  \providecommand{\bm}[1]{\symbfit{#1}}
\else
  \providecommand{\bm}[1]{\boldsymbol{#1}}
\fi

\providecommand{\oiint}{\oint\mkern-13mu\oint}
\providecommand{\ointctrclockwise}{\oint}

\DeclarePairedDelimiter{\abs}{\lvert}{\rvert}
\DeclarePairedDelimiter{\norm}{\lVert}{\rVert}
\DeclarePairedDelimiter{\innerp}{\langle}{\rangle}
\DeclarePairedDelimiter{\ceil}{\lceil}{\rceil}
\DeclarePairedDelimiter{\floor}{\lfloor}{\rfloor}

\DeclareMathOperator*{\argmin}{arg\,min}
\DeclareMathOperator*{\argmax}{arg\,max}
\DeclareMathOperator*{\Arg}{Arg}
\DeclareMathOperator*{\Real}{Re}
\DeclareMathOperator*{\Image}{Img}

\newcommand{\vv}[1]{\overrightarrow{#1}}
\newcommand{\vs}[1]{\vec{#1}}

%% ------------------------------------------------------------- 4. Cleveref & Cross-referencing
\makeatletter
\expandafter\let\csname opt@hyperref.sty\endcsname\@empty
\makeatother
\RequirePackage{cleveref}

\makeatletter
\let\utb@orig@href\href
\renewcommand{\href}[2]{\utb@orig@href{#1}{\textcolor{utbmain}{#2}}}

\renewcommand{\eqref}[1]{\hyperref[#1]{\textcolor{utbmain}{\textup{(\ref*{#1})}}}}

\crefname{equation}{phương trình}{phương trình}
\Crefname{equation}{Phương trình}{Phương trình}
\crefformat{equation}{#2\textcolor{utbmain}{phương trình~#1}#3}
\Crefformat{equation}{#2\textcolor{utbmain}{Phương trình~#1}#3}

\crefname{figure}{hình}{hình}
\Crefname{figure}{Hình}{Hình}
\crefformat{figure}{#2\textcolor{utbmain}{hình~#1}#3}
\Crefformat{figure}{#2\textcolor{utbmain}{Hình~#1}#3}

\crefname{table}{bảng}{bảng}
\Crefname{table}{Bảng}{Bảng}
\crefformat{table}{#2\textcolor{utbmain}{bảng~#1}#3}
\Crefformat{table}{#2\textcolor{utbmain}{Bảng~#1}#3}

\crefname{section}{mục}{mục}
\Crefname{section}{Mục}{Mục}
\crefformat{section}{#2\textcolor{utbmain}{mục~#1}#3}
\Crefformat{section}{#2\textcolor{utbmain}{Mục~#1}#3}

\crefname{subsection}{mục}{mục}
\Crefname{subsection}{Mục}{Mục}
\crefformat{subsection}{#2\textcolor{utbmain}{mục~#1}#3}
\Crefformat{subsection}{#2\textcolor{utbmain}{Mục~#1}#3}

\crefname{utbsqlisti}{mục}{mục}
\Crefname{utbsqlisti}{Mục}{Mục}
\crefformat{utbsqlisti}{#2\textcolor{utbmain}{mục~#1}#3}
\Crefformat{utbsqlisti}{#2\textcolor{utbmain}{Mục~#1}#3}

\crefname{hustsqlisti}{mục}{mục}
\Crefname{hustsqlisti}{Mục}{Mục}
\crefformat{hustsqlisti}{#2\textcolor{utbmain}{mục~#1}#3}
\Crefformat{hustsqlisti}{#2\textcolor{utbmain}{Mục~#1}#3}

\crefname{algocf}{thuật toán}{thuật toán}
\Crefname{algocf}{Thuật toán}{Thuật toán}
\crefformat{algocf}{#2\textcolor{utbmain}{thuật toán~#1}#3}
\Crefformat{algocf}{#2\textcolor{utbmain}{Thuật toán~#1}#3}

\crefname{lstlisting}{đoạn mã}{đoạn mã}
\Crefname{lstlisting}{Đoạn mã}{Đoạn mã}
\crefformat{lstlisting}{#2\textcolor{utbmain}{đoạn mã~#1}#3}
\Crefformat{lstlisting}{#2\textcolor{utbmain}{Đoạn mã~#1}#3}

\crefname{listing}{đoạn mã}{đoạn mã}
\Crefname{listing}{Đoạn mã}{Đoạn mã}
\crefformat{listing}{#2\textcolor{utbmain}{đoạn mã~#1}#3}
\Crefformat{listing}{#2\textcolor{utbmain}{Đoạn mã~#1}#3}

\crefname{theorem}{định lý}{định lý}
\Crefname{theorem}{Định lý}{Định lý}
\crefformat{theorem}{#2\textcolor{utbmain}{định lý~#1}#3}
\Crefformat{theorem}{#2\textcolor{utbmain}{Định lý~#1}#3}

\crefname{definition}{định nghĩa}{định nghĩa}
\Crefname{definition}{Định nghĩa}{Định nghĩa}
\crefformat{definition}{#2\textcolor{utbmain}{định nghĩa~#1}#3}
\Crefformat{definition}{#2\textcolor{utbmain}{Định nghĩa~#1}#3}

\crefname{lemma}{bổ đề}{bổ đề}
\Crefname{lemma}{Bổ đề}{Bổ đề}
\crefformat{lemma}{#2\textcolor{utbmain}{bổ đề~#1}#3}
\Crefformat{lemma}{#2\textcolor{utbmain}{Bổ đề~#1}#3}

\crefname{corollary}{hệ quả}{hệ quả}
\Crefname{corollary}{Hệ quả}{Hệ quả}
\crefformat{corollary}{#2\textcolor{utbmain}{hệ quả~#1}#3}
\Crefformat{corollary}{#2\textcolor{utbmain}{Hệ quả~#1}#3}

\crefname{example}{ví dụ}{ví dụ}
\Crefname{example}{Ví dụ}{Ví dụ}
\crefformat{example}{#2\textcolor{utbmain}{ví dụ~#1}#3}
\Crefformat{example}{#2\textcolor{utbmain}{Ví dụ~#1}#3}

%% 4.1. Universal Figure / Table Custom Caption & Cross-Reference Engine
\newcommand{\utbfigurename}{Hình}
\newcommand{\utbtablename}{Bảng}
\newcommand{\utbsetfigurename}[1]{\renewcommand{\utbfigurename}{#1}}
\newcommand{\utbsettablename}[1]{\renewcommand{\utbtablename}{#1}}

\newcounter{utbref}
\crefname{utbref}{}{}
\Crefname{utbref}{}{}
\crefformat{utbref}{#2\textcolor{utbmain}{#1}#3}
\Crefformat{utbref}{#2\textcolor{utbmain}{#1}#3}

\NewDocumentCommand{\utbcaption}{ O{\utbfigurename} o m }{%
  \par\vspace{3pt}\noindent
  \phantomsection
  \edef\utb@cap@name{#1}%
  \IfValueTF{#2}{%
    \edef\utb@cap@num{#2}%
  }{%
    \@ifundefined{c@utb@cnt@\utb@cap@name}{%
      \newcounter{utb@cnt@\utb@cap@name}%
    }{}%
    \stepcounter{utb@cnt@\utb@cap@name}%
    \edef\utb@cap@num{\arabic{utb@cnt@\utb@cap@name}}%
  }%
  \def\theutbref{\utb@cap@name~\utb@cap@num}%
  \refstepcounter{utbref}%
  \begin{center}
    {\scriptsize\bfseries\color{utbmain}\utb@cap@name~\utb@cap@num:}\ %
    {\normalfont\scriptsize #3}%
  \end{center}%
  \par\vspace{2pt}%
}

\let\hustcaption\utbcaption
\let\hustfigurename\utbfigurename
\let\husttablename\utbtablename
\let\hustsetfigurename\utbsetfigurename
\let\hustsettablename\utbsettablename
\makeatother

%% ------------------------------------------------------------- 5. Code & Pseudocode
\RequirePackage{listings}
\RequirePackage[linesnumbered,ruled,vlined]{algorithm2e}
\RequirePackage[highlightmode=immediate]{minted}

\SetNlSkip{0.5em}
\AtBeginEnvironment{algorithm}{%
  \IncMargin{1.2em}%
  \SetKwInput{KwIn}{\kern-1.2em Input}%
  \SetKwInput{KwOut}{\kern-1.2em Output}%
}

\makeatletter
\AtBeginDocument{%
  \renewcommand{\fnum@lstlisting}{Listing~\thelstlisting}%
}
\makeatother

\colorlet{utbcodebg}{utbmain!5!white}
\colorlet{hustcodebg}{utbcodebg}
\lstdefinestyle{utbcode}{
  backgroundcolor=\color{utbcodebg},
  basicstyle=\ttfamily\scriptsize,
  keywordstyle=\bfseries\color{utbmain},
  stringstyle=\color{UTBOrange},
  commentstyle=\itshape\color{UTBGray},
  numberstyle=\tiny\color{UTBGray},
  numbers=left,
  numbersep=5pt,
  framesep=3pt,
  xleftmargin=1.5em,
  framexleftmargin=\dimexpr 1.5em - 3.5pt\relax,
  xrightmargin=3.5pt,
  framexrightmargin=0pt,
  frame=single,
  framerule=0.5pt,
  rulecolor=\color{utbmain!30},
  breaklines=true,
  tabsize=2,
  showstringspaces=false
}
\lstdefinestyle{hustcode}{style=utbcode}

\setminted{
  fontsize=\scriptsize,
  linenos=true,
  numbersep=5pt,
  xleftmargin=1.2em,
  breaklines=true,
  tabsize=2,
  frame=lines,
  framerule=0.5pt,
  rulecolor=\color{utbmain!40}
}
\renewcommand{\theFancyVerbLine}{\textcolor{UTBGray}{\tiny\arabic{FancyVerbLine}}}

%% ------------------------------------------------------------- 6. Beamer Theorems & Academic Blocks
\setbeamertemplate{theorems}[numbered]

\deftranslation{Definition}{Định nghĩa}
\deftranslation{Theorem}{Định lý}
\deftranslation{Lemma}{Bổ đề}
\deftranslation{Corollary}{Hệ quả}
\deftranslation{Example}{Ví dụ}
\deftranslation{Examples}{Ví dụ}
\deftranslation{Proof}{Chứng minh}
\deftranslation[to=English]{Definition}{Định nghĩa}
\deftranslation[to=English]{Theorem}{Định lý}
\deftranslation[to=English]{Lemma}{Bổ đề}
\deftranslation[to=English]{Corollary}{Hệ quả}
\deftranslation[to=English]{Example}{Ví dụ}
\deftranslation[to=English]{Examples}{Ví dụ}
\deftranslation[to=English]{Proof}{Chứng minh}
\providecommand{\proofname}{Chứng minh}
\renewcommand{\proofname}{Chứng minh}

%% Vector warning icon:
\DeclareRobustCommand{\warning}{%
  \tikz[baseline=0pt]{%
    \draw[line width=0.6pt, line join=round, draw=UTBRed, fill=UTBYellow!70, inner sep=0pt]
      (0,0) -- (1.1em,0) -- (0.55em,0.95em) -- cycle;
    \node[font=\tiny\bfseries, text=UTBRed] at (0.55em,0.3em) {!};
  }\xspace}

%% End of utb-setup.tex


\usepackage[backend=biber,style=numeric,sorting=none,maxcitenames=2,mincitenames=1]{biblatex}
\addbibresource{refs.bib}

\title[Xây dựng mẫu trình chiếu Beamer chuẩn nhận diện UTB]{{\normalsize ĐỒ ÁN TỐT NGHIỆP CỬ NHÂN}\\[5pt] {\bfseries Xây dựng mẫu trình chiếu Beamer chuẩn nhận diện UTB}}
\subtitle{Hệ thống trình chiếu LaTeX hiện đại, linh hoạt tỉ lệ và kích thước phông chữ}
\author[Lò Mạnh Đạt]{Lò Mạnh Đạt - 20231234 \\ GVHD: PGS. TS. Lê Văn Thu C}
\institute{Khoa Tự nhiên và Công nghệ,\\Trường Đại học Tây Bắc}
\date{\today}

\begin{document}

%% ------------------------------------------------------------- 0a. Welcome Screen 1: Normal bg-dark
\begin{utbwelcomeframe}
\end{utbwelcomeframe}

%% ------------------------------------------------------------- 0b. Welcome Screen 2: Modified bg-background showcase
\begin{utbwelcomeframe}[bg=background, footer={DEMO SHOWCASE}]
\end{utbwelcomeframe}

%% ------------------------------------------------------------- 1. Title Slide (Image 1)
\utbsetup{slogantext={}}
\begin{frame}[plain,noframenumbering]
  \titlepage
\end{frame}

%% ------------------------------------------------------------- 2. Outline (No Progress Bar on TOC)
\begin{frame}[noprogressbar]{\utbtoctitle}{\utbtocsubtitle}
  \tableofcontents[split=3,hideallsubsections]
\end{frame}

%% ------------------------------------------------------------- Section 1
\section[Quy chuẩn]{Giới thiệu \& quy chuẩn}
\subsection{Hệ thống nhận diện thương hiệu}
\begin{frame}{Tạo cụm logo động theo sổ tay thương hiệu}{Quy chuẩn Trường Đại học Tây Bắc (Việt - Anh) và Đơn vị thành viên}
  \scriptsize
  \begin{columns}[T]
    \begin{column}{0.48\textwidth}
      \textbf{1. Cụm chính thức Việt - Anh (1 dòng):}
      \vskip0.3em
      \utblogo[height=0.85cm]
      \vskip0.7em
      \textbf{2. Bản 2 dòng Việt - Anh:}
      \vskip0.3em
      \utblogo[height=0.85cm, layout=2, vn={TRƯỜNG ĐẠI HỌC\newline TÂY BẮC}, en={TAY BAC\newline UNIVERSITY}]
      \vskip0.7em
      \textbf{3. Đơn vị thành viên 1 dòng:}
      \vskip0.3em
      \utblogo[height=0.85cm, unit={KHOA KHOA HỌC TỰ NHIÊN - CÔNG NGHỆ}]
    \end{column}

    \begin{column}{0.48\textwidth}
      \textbf{4a. Đơn vị 2 dòng -- ĐHTB 2 dòng (Đối xứng):}
      \vskip0.3em
      \utblogo[height=0.85cm, layout=4a, unit={KHOA KHOA HỌC TỰ NHIÊN\newline VÀ CÔNG NGHỆ}]
      \vskip0.7em
      \textbf{4b. Đơn vị 2 dòng -- ĐHTB 1 dòng (Căn giữa):}
      \vskip0.3em
      \utblogo[height=0.85cm, layout=4b, unit={KHOA KHOA HỌC TỰ NHIÊN\newline VÀ CÔNG NGHỆ}]
      \vskip0.7em
      \textbf{5. Cụm logo Khoa KHTN - CN (với huy hiệu riêng):}
      \vskip0.3em
      \utblogo[image=nst-shield-utb.png, height=0.85cm, vn={KHOA KHOA HỌC TỰ NHIÊN - CÔNG NGHỆ}, en={Faculty of Natural Sciences and Technology}, color=utbnavy, equalwidth=false]
    \end{column}
  \end{columns}
\end{frame}

\subsection{Hệ thống màu sắc nhận diện thương hiệu}
\begin{frame}{Hệ thống màu sắc nhận diện thương hiệu UTB}{Bộ mã màu quy chuẩn theo Hệ thống nhận diện Trường Đại học Tây Bắc}
  \begin{columns}[T]
    \begin{column}{0.50\textwidth}
      \textbf{Màu sắc chủ đạo \& thứ cấp}:
      \begin{itemize}\setlength{\itemsep}{3pt}
        \item \textcolor{UTBGreen}{$\blacksquare$ \texttt{UTBGreen} / \texttt{utbmain}} (\texttt{\#12711F})\\
          \scriptsize Xanh lá đậm nhận diện cốt lõi của Trường Đại học Tây Bắc.
        \item \textcolor{UTBNavy}{$\blacksquare$ \texttt{UTBNavy}} (\texttt{\#0B2D4F})\\
          \scriptsize Xanh navy màu nhận diện thứ hai, biểu trưng tri thức và uy tín.
        \item \textcolor{UTBDarkGreen}{$\blacksquare$ \texttt{UTBDarkGreen}} (\texttt{\#105419})\\
          \scriptsize Xanh đậm nền banner dàn ý slide và chuyển tiếp.
        \item \textcolor{UTBLightGreen}{$\blacksquare$ \texttt{UTBLightGreen}} (\texttt{\#CDFAB1})\\
          \scriptsize Xanh non màu nhạt thanh tiến trình và điểm nhấn.
      \end{itemize}
    \end{column}

    \begin{column}{0.48\textwidth}
      \textbf{Màu trung tính \& sắc thái bổ trợ}:
      \begin{itemize}\setlength{\itemsep}{3pt}
        \item \textcolor{UTBBlack}{$\blacksquare$ \texttt{UTBBlack}} (\texttt{\#231F20})\\
          \scriptsize Đen nhận diện tiêu chuẩn thay thế đen thuần 100\%.
        \item \textcolor{UTBWhite}{$\blacksquare$ \texttt{UTBWhite}} (\texttt{\#FFFFFF})\\
          \scriptsize Màu trắng chuẩn phông nền và độ tương phản cao.
        \item \textcolor{UTBOrange}{$\blacksquare$ \texttt{UTBOrange}} (\texttt{\#F57F20})\\
          \scriptsize Cam năng động bổ trợ cho các khối ví dụ minh họa.
        \item \textcolor{UTBGray}{$\blacksquare$ \texttt{UTBGray}} (\texttt{\#A6A6A6})\\
          \scriptsize Xám bổ trợ đường phân cách và chú thích.
      \end{itemize}
    \end{column}
  \end{columns}
\end{frame}

%% ------------------------------------------------------------- Section 2
\section[Giải thuật]{Các khối nội dung \& giải thuật}
\subsection{Các loại khối hiển thị}
\begin{frame}{Các loại khối hiển thị (Beamer blocks)}{Đồng bộ màu sắc nhận diện thương hiệu UTB}
  \begin{block}{Khối tiêu chuẩn (utbmain)}
    Khối định nghĩa hoặc thông tin quan trọng được làm nổi bật với màu xanh lá nhận diện của Trường Đại học Tây Bắc.
  \end{block}

  \begin{alertblock}{Khối cảnh báo (utbalert)}
    Khối cảnh báo hoặc kết luận then chốt sử dụng màu xanh navy nhận diện thứ hai của UTB.
  \end{alertblock}

  \begin{exampleblock}{Khối ví dụ (utbexample)}
    Ví dụ minh họa, mã nguồn hoặc kết quả thực nghiệm với màu cam bổ trợ UTB.
  \end{exampleblock}
\end{frame}

\subsection{Môi trường định nghĩa \& khối ghi chú}
\begin{frame}{Môi trường định nghĩa và khối ghi chú}{Đồng bộ màu sắc nhận diện và chú thích tự động}
  \begin{definition}[Phép toán hai ngôi]
    Cho tập $X \neq \varnothing$. Một \note{phép toán hai ngôi}{Một ánh xạ từ $X^2 \to X$ gán mỗi cặp phần tử $(x, y)$ thành một phần tử $x * y \in X$.} trên tập $X$, ký hiệu là $*$, là một quy tắc biến đổi mỗi cặp phần tử $(x, y) \in X^2$ thành phần tử $x * y \in X$. Nói cách khác, phép toán $*$ là một ánh xạ:
    \[
      (*) : X^2 \longrightarrow X, \quad (x, y) \longmapsto x * y
    \]
  \end{definition}

  \begin{exampleblock}{Nhận xét quan trọng}
    \warning Các môi trường khối và toán học chuẩn của Beamer (\texttt{block}, \texttt{definition}, \texttt{theorem}, \texttt{exampleblock}) được tinh chỉnh tự động theo bảng màu thương hiệu UTB, hỗ trợ đánh số và tham chiếu chéo thông minh qua \texttt{cleveref}.
  \end{exampleblock}
\end{frame}

\subsection{Môi trường toán học chuẩn Beamer}
\begin{frame}{Môi trường toán học chuẩn Beamer}{Bộ môi trường định lý, bổ đề, ví dụ và chứng minh đồng bộ}
  \scriptsize
  \begin{columns}[T]
    \begin{column}{0.48\textwidth}
      \begin{theorem}[Hội tụ nghiệm Pareto]
        Dãy lặp $\{x_k\}$ sinh bởi thuật toán MPG-E hội tụ yếu về nghiệm tối ưu $x^* \in C$.
      \end{theorem}
      \vskip 2pt
      \begin{exampleblock}{Tính đơn điệu}
        Hàm khoảng cách $\|x_k - x^*\|$ giảm đơn điệu theo chỉ số bước lặp $k$.
      \end{exampleblock}
    \end{column}

    \begin{column}{0.48\textwidth}
      \begin{lemma}[Bất đẳng thức chiếu]
        Với mọi $z \in \mathcal{X}$, ta có $\innerp{x - \mathrm{Proj}_{\mathcal{X}}(x)}{\mathrm{Proj}_{\mathcal{X}}(x) - z} \ge 0$.
      \end{lemma}
      \vskip 2pt
      \begin{proof}
        Áp dụng tính chất toán tử chiếu không giãn trên tập lồi đóng trong không gian Hilbert.
      \end{proof}
    \end{column}
  \end{columns}
\end{frame}

\subsection{Tùy biến danh sách (enumitem \& tasks)}
\begin{frame}{Tùy biến danh sách với enumitem và tasks}{Đánh số dạng ô vuông nhận diện thương hiệu UTB}
  \begin{definition}[Nghiệm Pareto]
    \begin{utbsqlist}
      \item $x^* \in C$ là \emph{nghiệm Pareto} nếu $\nexists x \in C$ sao cho $g_i(x) \le g_i(x^*)$ với mọi $i$ và $g_j(x) < g_j(x^*)$ tại ít nhất một chỉ số $j$. \label{def:pareto}
      \item $\bar{x} \in C$ là \emph{nghiệm Pareto yếu} nếu $\nexists x \in C$ sao cho $g_i(x) < g_i(\bar{x})$ với mọi $i$.
      \item Tập hợp tất cả các nghiệm Pareto được gọi là \emph{tập Pareto}. Ảnh của nó trong không gian mục tiêu được gọi là \emph{biên Pareto}.
    \end{utbsqlist}
  \end{definition}
  \vskip 2pt
  \textbf{Phân loại bài toán tối ưu với \texttt{tasks}} (gọi tùy biến với badge ô vuông chữ cái):
  \begin{tasks}[label=\utbitemsq{\alph*}, label-width=18pt, item-indent=22pt, label-format={\bfseries\color{white}}](2)
    \task Bài toán lồi đa mục tiêu
    \task Bài toán không lồi với gradient
    \task Điều kiện tối ưu Karush-Kuhn-Tucker
    \task Thuật toán chiếu điểm gần kề
  \end{tasks}
\end{frame}

\subsection{Toán học chuyên sâu \& biểu đồ TikZ}
\begin{frame}{Toán học chuyên sâu, tích phân đường cong kín \& TikZ}{Công thức giải tích phức, đạo hàm và tham chiếu chéo}
  \begin{columns}[c]
    \begin{column}{0.62\textwidth}
      Áp dụng cho hàm giải tích $f(z)$ trên miền đơn liên có biên cong kín $C$:
      \begin{equation}
        \ointctrclockwise_{C} \frac{f(z)}{z - z_0}\,\mathrm{d}z = 2\pi i f(z_0) \label{eq:cauchy}
      \end{equation}
      \begin{itemize}\setlength{\itemsep}{2pt}
        \item Ký hiệu vector in đậm $\bm{w} \in \mathbb{C}^n$, đạo hàm $\odv{f}{z}$, đạo hàm riêng $\pdv{u}{x} = \pdv{v}{y}$.
        \item Chuẩn $\norm{z}$, tích vô hướng $\innerp{u}{v}$, tích phân kép $\oiint_{\mathcal{S}} \bm{F}\cdot\mathrm{d}\bm{S}$.
        \item Tham chiếu chéo tự động: \Cref{eq:cauchy} (hoặc \eqref{eq:cauchy}), \cref{def:pareto}, và \cref{fig:cauchy}.
        \item Liên kết ngoài trực tiếp: \href{https://utb.edu.vn}{\underline{Trang chủ Trường Đại học Tây Bắc}}.
      \end{itemize}
    \end{column}
    \begin{column}{0.36\textwidth}
      \centering
      \begin{tikzpicture}[scale=0.95, baseline=0]
        \draw[->, >=Stealth, UTBGray] (-1.8,0) -- (1.8,0) node[right, font=\tiny] {$\mathrm{Re}$};
        \draw[->, >=Stealth, UTBGray] (0,-1.4) -- (0,1.4) node[above, font=\tiny] {$\mathrm{Im}$};
        \draw[thick, utbmain, postaction={decorate, decoration={markings, mark=at position 0.35 with {\arrow{Stealth}}}}]
          (0,0) circle (1.0cm);
        \fill[utbalert] (0.2, 0.3) circle (1.5pt) node[above right, font=\scriptsize] {$z_0$};
        \node[utbmain, font=\scriptsize] at (-0.8, 0.8) {$C$};
      \end{tikzpicture}
      \vskip 2pt
      \utbcaption[Hình][1]{Miền giải tích $C$}\label{fig:cauchy}
    \end{column}
  \end{columns}
\end{frame}

\subsection{Bảng dữ liệu nâng cao}
\begin{frame}{Bảng dữ liệu nâng cao với booktabs và tabularx}{Tự động dãn độ rộng cột, gộp ô hàng ngang và hàng dọc}
  \small
  So sánh hiệu năng thực nghiệm theo các nghiên cứu của Alvarez\cite{alvarez2024} và Bello-Cruz\cite{bellocruz2025}:
  \vskip 0.3em
  \centering
  \begin{tabularx}{\textwidth}{l c c X r}
    \toprule
    \textbf{Phương pháp} & \textbf{Số vòng lặp} & \textbf{Thời gian (s)} & \textbf{Ghi chú thuật toán} & \textbf{Độ chính xác} \\
    \midrule
    \multirow{2}{*}{Baseline} & 150 & 12.4 & Gradient Descent cổ điển & 87.5\% \\
                              & 120 & 10.1 & Thêm điều kiện bước nhảy Armijo & 89.2\% \\
    \midrule
    \multirow{2}{*}{Quán tính} & 85  & 6.8  & Nesterov Accelerated Gradient & 93.4\% \\
                               & 60  & 4.5  & Heavy-ball trong không gian Hilbert\cite{alvarez2024} & \textbf{95.1\%} \\
    \midrule
    \rowcolor{utbmain!10}
    \textbf{MPG-E (Đề xuất)} & \textbf{38} & \textbf{2.9} & Chiếu gradient đa mục tiêu\cite{bellocruz2025} & \textbf{96.8\%} \\
    \bottomrule
  \end{tabularx}
  \vskip 2pt
  \utbcaption[Bảng][2.5]{So sánh hiệu năng thực nghiệm thuật toán}\label{tab:exp}
\end{frame}

\subsection{Mã giải thuật \& lập trình}
\begin{frame}[fragile]{Biểu diễn giải thuật tối ưu hóa đa mục tiêu với algorithm2e, listings và minted trong trình chiếu khoa học}{Cấu trúc mã giả khoa học và khối mã nguồn tô màu cú pháp đa ngôn ngữ}
  \begin{columns}[T]
    \begin{column}{0.49\textwidth}
      \begin{minipage}{\linewidth}
      \begin{algorithm}[H]
        \scriptsize
        \caption{Thuật toán MPG-E}\label{algo:mpg}
        \KwIn{$x_0 \in \mathcal{X}$, $\beta \in (0,1)$, $\gamma > 0$}
        \KwOut{Nghiệm Pareto yếu $x^*$}
        \BlankLine
        \For{$k = 0, 1, 2, \dots$}{
          $d_k = \argmin_{d} \left\{ \max_i \innerp{\nabla f_i(x_k)}{d} + \frac{1}{2}\norm{d}^2 \right\}$\;
          \If{$\norm{d_k} < \epsilon$}{
            \Return $x_k$\;
          }
          $x_{k+1} = \mathrm{Proj}_{\mathcal{X}}(x_k + \alpha_k d_k)$\;
        }
      \end{algorithm}
      \end{minipage}
    \end{column}

    \begin{column}{0.49\textwidth}
\begin{lstlisting}[style=utbcode, language=Python]
import numpy as np
def mpg_step(x, grad, alpha=0.01):
    d_k = -np.mean(grad(x), axis=0)
    return x + alpha * d_k
\end{lstlisting}
\vskip -4pt
\utbcaption[Listing][1]{Hàm mục tiêu Python (\texttt{listings})}\label{lst:python}
\vskip 2pt
\begin{minted}{cpp}
double clip_box(double x, double r) {
    // Chieu len mien [-r, r]
    return std::clamp(x, -r, r);
}
\end{minted}
\vskip -4pt
\utbcaption[Listing][2]{Hàm C++ (\texttt{minted}, cần Python \& \texttt{-shell-escape})}\label{lst:cpp}
    \end{column}
  \end{columns}
\end{frame}

%% ------------------------------------------------------------- Section 3: Biber / BibLaTeX Citations
\section{Trích dẫn \& tài liệu tham khảo}
\subsection{Bố cục wrapfig \& trích dẫn}
\begin{frame}{Trích dẫn tài liệu với bố cục wrapfig}{Bọc văn bản xung quanh logo với số dòng quy định}
  \begin{wrapfigure}[8]{r}{0.16\textwidth}
    \vspace{-1.2em}
    \centering
    \includegraphics[width=0.13\textwidth]{assets/utb-shield.png}
  \end{wrapfigure}
  Nghiên cứu về phương pháp tối ưu hóa đa mục tiêu được đề xuất bởi Bello-Cruz và cộng sự\cite{bellocruz2025} nhằm giải quyết bài toán tìm điểm hữu hiệu trong không gian Hilbert. Mẫu trình chiếu Beamer được phát triển tại \note{Đại học Tây Bắc}{Trường Đại học công lập đa ngành tại khu vực Tây Bắc, thành lập 1960}\cite{utb2026} với hệ thống nhận diện đồng bộ.

  Hình ảnh huy hiệu Trường Đại học Tây Bắc ở bên phải được bao bọc bởi chính xác 6 dòng văn bản nhờ gói \texttt{wrapfig}, đảm bảo sự hài hòa về mặt thị giác trên khung hình thuyết trình.
\end{frame}

\subsection{Phương pháp tối ưu hóa đa mục tiêu}
\begin{frame}{Nghiên cứu tối ưu hóa đa mục tiêu}{Hệ thống tham chiếu chéo toàn diện và trích dẫn tự động}
  Xét bài toán tối ưu hóa đa mục tiêu với tập ràng buộc $C \subset \mathbb{R}^n$:
  \begin{equation}
    \min_{x \in C} \; F(x) = \big(f_1(x), f_2(x), \dots, f_m(x)\big) \label{eq:mop}
  \end{equation}
  Nghiên cứu của Bello-Cruz và cộng sự\cite{bellocruz2025} giải quyết \Cref{eq:mop} theo tiêu chuẩn \cref{def:pareto}:
  \begin{itemize}\setlength{\itemsep}{3pt}
    \item \textbf{Mô hình \& hình học}: Bài toán \Cref{eq:mop} liên kết với miền giải tích \cref{fig:cauchy} qua công thức tích phân vi phân \eqref{eq:cauchy}.
    \item \textbf{Thuật toán \& cài đặt}: Quy trình giải lặp quán tính\cite{alvarez2024} được đặc tả tại \cref{algo:mpg}, mã nguồn mẫu tại \cref{lst:python} và \cref{lst:cpp}.
    \item \textbf{Hiệu năng thực nghiệm}: Bảng so sánh \cref{tab:exp} khẳng định thuật toán hội tụ về nghiệm Pareto \cref{def:pareto} trên hệ thống chuẩn\cite{utb2026}.
    \item \textbf{Chân trang thông minh}: Cụm trích dẫn tránh logo (\texttt{footerlogo=mark}), tự động rút gọn \texttt{...} trên 1 dòng.
  \end{itemize}
\end{frame}

\subsection{Thuật toán quán tính trong không gian Hilbert}
\begin{frame}{Nhiều trích dẫn trên cùng một slide}{Khóa tối đa 4 dòng chân trang: tự động chia 2 cột và rút gọn \texttt{...} khi có $>4$ mục}
  Các thuật toán giải bài toán tối ưu với định dạng \texttt{tasks} mặc định:
  \begin{tasks}(2)
    \task Phương pháp gradient gần kề\cite{bellocruz2025}
    \task Thuật toán quán tính Hilbert\cite{alvarez2024}
    \task Chuẩn hóa nhận diện Trường Đại học Tây Bắc\cite{utb2026}
    \task Phương pháp đối ngẫu tăng tốc\cite{nesterov2018}
  \end{tasks}
  \vskip 0.5em
  \begin{itemize}
    \item Môi trường \texttt{tasks} mặc định kế thừa dấu đầu dòng hình vuông đặc trưng của Beamer itemize.
    \begin{itemize}[nosep]
      \item Hộp chân trang khóa tối đa \textbf{4 dòng}: tự động rút gọn tên bài báo dài bằng \texttt{...} và chuyển sang \note{bố cục 2 cột}{Tự động kích hoạt khi số trích dẫn và ghi chú trên slide vượt quá 4 dòng} (tối đa 8 mục).
      \item Phân cấp danh sách cấp 3 (dùng \texttt{enumerate} để \note{đánh số La Mã}{Bộ đếm cấp 3 tự động tăng i, ii, iii chính xác} tự động tăng chính xác):
      \begin{enumerate}[label = \roman*)]
        \item Bước kiểm thử thứ nhất (\texttt{i})
        \item Bước kiểm thử kế tiếp tự động tăng đúng số (\texttt{ii})
      \end{enumerate}
    \end{itemize}
  \end{itemize}
\end{frame}

%% ------------------------------------------------------------- Section 4 (Tắt panel slide để minh họa tùy chọn)
\utbsection[panel=false, alias=Tùy biến]
\section{Tùy biến \& linh hoạt}

\subsection{Bảng tùy chọn theme UTB}
\begin{frame}{Các tùy chọn của gói theme}{Tùy biến dễ dàng qua \texttt{\textbackslash usetheme} hoặc \texttt{\textbackslash utbsetup}}
  \scriptsize
  \centering
  \begin{tabularx}{\textwidth}{>{\ttfamily\bfseries\color{utbmain}}r >{\small}X}
    color       & \texttt{green} (mặc định) hoặc \texttt{navy}. \\
    font        & \texttt{rokkit} (mặc định), \texttt{lato}, \texttt{cm}, \texttt{times}, hoặc \texttt{arial}. \\
    titlefont   & \texttt{avantgarde} (mặc định theo sổ tay) hoặc \texttt{same}. \\
    titlepage   & \texttt{cover} (nền chấm đậm), \texttt{light} (nền chấm mờ cho 4:3), \texttt{solid} (nền màu đặc), \texttt{dark}, \texttt{welcome}. \\
    bar         & \texttt{thin} (thanh mảnh) hoặc \texttt{thick} (thanh dày cho phụ đề). \\
    progressbar & \texttt{on} (bật thanh tiến trình) hoặc \texttt{off}. \\
    sectionpage & \texttt{panel} (tạo trang chuyển mục) hoặc \texttt{none}. \\
    footermeta  & \texttt{on} (hiện tác giả, tiêu đề, ngày ngắn ở chân trang) hoặc \texttt{off}. \\
    headerright & \texttt{off}, \texttt{pagenumber} (số trang lên góc phải tiêu đề), \texttt{logo} (\texttt{headerlogo=...}). \\
    framebg     & \texttt{on} (họa tiết thoi chéo chìm cho frame \& panel) hoặc \texttt{off}. \\
    slogan      & \texttt{on} (hiển thị TRƯỜNG ĐẠI HỌC TÂY BẮC) hoặc \texttt{off}. \\
    footerlogo  & \texttt{mark} (mặc định), \texttt{full}, \texttt{wordmark}, \texttt{utb}, hoặc \texttt{none}. \\
    numbering   & \texttt{fraction} ($n/N$), \texttt{plain} ($n$), hoặc \texttt{none}. \\
    logolayout  & \texttt{1} (1 dòng), \texttt{2} (2 dòng), \texttt{3} (đơn vị 1 dòng), \texttt{4a}, \texttt{4b}. \\
    unit        & Tên khoa/viện thành viên (1 dòng hoặc 2 dòng kèm \texttt{\textbackslash newline}). \\
    dateaddress & \texttt{on} (kèm địa danh tiếng Việt) hoặc \texttt{off}. \\
    address     & Địa danh đi kèm ngày tháng (ví dụ: \texttt{\{Sơn La\}}). \\
  \end{tabularx}
\end{frame}

\subsection{Các kiểu khung hình (Frame Types)}
\begin{frame}{Các kiểu khung hình và môi trường slide}{Danh mục khung hình chuyên dụng kèm tùy chọn đối số}
  \scriptsize
  \centering
  \begin{tabularx}{\textwidth}{>{\ttfamily\bfseries\color{utbmain}}r >{\small}X}
    frame               & Khung cơ bản: \texttt{[t|c|b]}, \texttt{[plain]}, \texttt{[fragile]}, \texttt{[noframenumbering]}, \texttt{[noprogressbar]}. \\
    utbwelcomeframe    & Khung mở đầu / chào mừng: \texttt{[bg=dark|background|cover]}, \texttt{[footer=true|false]}. \\
    utbdarkframe       & Khung nền màu đậm nhận diện thương hiệu (slide kết luận, Q\&A). \\
    utbbackgroundframe & Khung nền hoa văn chìm dải sóng đồ họa UTB toàn trang. \\
    utbreferencesframe & Khung tài liệu tham khảo phân trang thủ công: \texttt{[from=N, to=M]}. \\
    utbsection         & Điều khiển bật/tắt slide chuyển mục cho \texttt{\textbackslash section}: \texttt{[panel=true|false]}. \\
    \textbackslash note & Ghi chú chân trang chữ cái ($a, b, \dots$) đồng bộ vị trí cùng \texttt{\textbackslash cite}. \\
  \end{tabularx}
\end{frame}

\subsection{Gợi ý phối màu \& tỉ lệ Beamer}
\begin{frame}{Gợi ý phối màu và tỉ lệ hiển thị Beamer}{Hệ màu bổ trợ mở rộng, phông chữ và các tỉ lệ màn hình hỗ trợ}
  \small
  \begin{columns}[T]
    \begin{column}{0.48\textwidth}
      \textbf{Màu nền gợi ý}:
      \begin{itemize}\setlength{\itemsep}{1pt}
        \item \colorbox{utbmain!8}{\texttt{utbmain!8}}: Nền khối tiêu chuẩn (\texttt{block}).
        \item \colorbox{utbalert!10}{\texttt{utbalert!10}}: Nền khối cảnh báo (\texttt{alertblock}).
        \item \colorbox{utbexample!12}{\texttt{utbexample!12}}: Nền khối ví dụ (\texttt{exampleblock}).
      \end{itemize}
      \vskip 0.3em
      \textbf{Khuyến nghị phông chữ}:
      \begin{itemize}\setlength{\itemsep}{1pt}
        \item \texttt{rokkit}: Phông nhận diện UTB, trang trọng và hiện đại.
        \item \texttt{lato}: Thanh thoát, bài giảng công nghệ.
        \item \texttt{cm}: Kinh điển, chuẩn công thức toán.
        \item \texttt{times}: Trang trọng, đồ án tốt nghiệp.
        \item \texttt{arial}: Nét đậm, dễ đọc từ khoảng cách xa.
      \end{itemize}
    \end{column}

    \begin{column}{0.48\textwidth}
      \textbf{Tất cả tỉ lệ Beamer hỗ trợ}:
      \begin{itemize}\setlength{\itemsep}{1pt}
        \item \texttt{169}: 16:9 widescreen màn hình LED, TV.
        \item \texttt{1610}: 16:10 màn hình laptop hiện đại.
        \item \texttt{149}: 14:9 tỉ lệ trung gian cân đối.
        \item \texttt{54}: 5:4 màn hình máy trạm chuyên dụng.
        \item \texttt{43}: 4:3 máy chiếu truyền thống.
        \item \texttt{32}: 3:2 máy tính bảng Surface.
        \item \texttt{141}: 1.41:1 chuẩn tương đương khổ giấy A4.
      \end{itemize}
    \end{column}
  \end{columns}
\end{frame}

\subsection{Yêu cầu cài đặt \& biên dịch}
\begin{frame}{Yêu cầu phần mềm và quy trình biên dịch}{Hướng dẫn thiết lập môi trường TeXstudio, MiKTeX và Overleaf}
  \begin{columns}[T]
    \begin{column}{0.50\textwidth}
      \textbf{1. Phần mềm \& Ứng dụng hỗ trợ}:
      \begin{itemize}\setlength{\itemsep}{2pt}
        \item \textbf{TeXstudio} (Khuyến nghị sử dụng offline).
        \item \textbf{MiKTeX / TeX Live} (Hệ phân phối LaTeX).
        \item \textbf{Overleaf}: Hoạt động trực tiếp trên trình duyệt.
      \end{itemize}
      \vskip 0.3em
      \textbf{2. Trình biên dịch bắt buộc}:
      \begin{itemize}\setlength{\itemsep}{2pt}
        \item \textbf{LuaLaTeX} (Khuyến nghị số 1) hoặc \textbf{XeLaTeX} để nạp phông hệ thống và tiếng Việt Unicode.
        \item Không dùng \texttt{pdfLaTeX} (không nạp được OTF/TTF).
      \end{itemize}
    \end{column}

    \begin{column}{0.48\textwidth}
      \textbf{3. Trình xử lý tài liệu tham khảo}:
      \begin{itemize}\setlength{\itemsep}{2pt}
        \item Sử dụng \textbf{Biber} kết hợp gói \texttt{biblatex}.
        \item Quy trình biên dịch:
          \[
            \text{Lua\LaTeX} \longrightarrow \text{Biber} \longrightarrow \text{Lua\LaTeX}
          \]
      \end{itemize}
      \vskip 0.3em
      \textbf{4. Cấu hình TeXstudio (1 chạm)}:
      \begin{itemize}\scriptsize\setlength{\itemsep}{2pt}
        \item \textit{Options $\to$ Configure TeXstudio $\to$ Build}:
        \item \textbf{Default Compiler}: chọn \texttt{LuaLaTeX} (thêm cờ \texttt{-shell-escape} và cài \textbf{Python} nếu dùng \texttt{minted}).
        \item \textbf{Default Bibliography}: chọn \texttt{Biber}.
        \item Nhấn \textbf{F5} để biên dịch toàn bộ.
      \end{itemize}
    \end{column}
  \end{columns}
\end{frame}

%% ------------------------------------------------------------- 5. References Frames (Manual Split with Identical Title)
\begin{utbreferencesframe}[from=1, to=2]
\end{utbreferencesframe}

\begin{utbreferencesframe}[from=3, to=4]
\end{utbreferencesframe}

%% ------------------------------------------------------------- Ending Dark Slide
\begin{utbdarkframe}
  \centering
  \vspace*{0.2\paperheight}
  {\Huge\bfseries\color{white} CẢM ƠN QUÝ THẦY CÔ\\VÀ CÁC BẠN ĐÃ LẮNG NGHE!\par}
  \vspace{1.5em}
  {\large\color{white!85} Q \& A\par}
\end{utbdarkframe}

\end{document}
